\documentclass[12pt]{book}
\usepackage[spanish]{babel}
\usepackage[top=3.5cm, bottom=3cm, left=2cm, right=2cm, headheight=2.5cm, footskip=1.5cm, marginparwidth=2cm]{geometry}
\usepackage{fancyhdr}
\usepackage{amsmath}
\usepackage{tcolorbox}
\usepackage{tikz}
\usepackage{amssymb}
\usepackage{fancybox}
\usepackage{lastpage}
\usepackage{xcolor}
\usepackage{tikz}
\usepackage{tikzpagenodes}
\usepackage{lipsum}
\usepackage{marginnote}
\usepackage{cancel}
\usepackage{multicol}
\usepackage{mdframed}
\usepackage{pgfplots}
\usepackage{hyperref}
\usepackage{etoolbox}
\usetikzlibrary{patterns} 
\tcbuselibrary{skins}
\usetikzlibrary{calc,angles,quotes}
\usepackage{tzplot}
\usepackage{float}
\usepackage{kinematikz}
\usepackage{titlesec}
\usepgfplotslibrary{fillbetween}
\usetikzlibrary{plotmarks}
\usepgfplotslibrary{polar}
\tcbuselibrary{theorems}

\definecolor{cadmiumred}{rgb}{0.89, 0.0, 0.13}
\definecolor{gre}{RGB}{101, 191, 127}
\definecolor{gree}{RGB}{7, 135, 44}
\definecolor{safetyorange(blazeorange)}{rgb}{1.0, 0.4, 0.0}
\definecolor{gold(web)(golden)}{rgb}{1.0, 0.84, 0.0}
\definecolor{bleudefrance}{rgb}{0.19, 0.55, 0.91}
\definecolor{kellygreen}{rgb}{0.3, 0.73, 0.09}


\newcommand{\ChapterStyle}[2]{%
  \begin{tikzpicture}[remember picture,overlay]
    % Fondo decorativo
    \fill[safetyorange(blazeorange)] (-3,0) rectangle (40,4); 
    \draw[blue, line width=2mm] (-3,4) -- (40,4);
    \draw[blue, line width=2mm] (-3,0) -- (40,0);
    \fill[white] (15,2) circle (1.7);
    \draw[line width=1mm] (15,2) circle (1.7)
    \draw[line width=0.5mm] (15,2) circle (1.5)
    \fill[kellygreen] (15,2) circle (1.5)
    \node at (15,2) {\Huge\bfseries\sffamily\color{white}{\thechapter}}
   \node[anchor=west] at (2,2) {\Huge\bfseries\sffamily\color{white}{#2}};
  \end{tikzpicture}
}

\titleformat{\chapter}[block]
  {\normalfont}
  {}
  {0pt}
  {%
    \ChapterStyle{\thechapter}{} % Pasa el número del capítulo y el título al diseño
  }


\newcommand{\Estilo}{
 \begin{tikzpicture}[remember picture, overlay]
         \draw
     (current page header area.south west) -- (current page header area.south east)
    \end{tikzpicture}

    \begin{tikzpicture}[remember picture, overlay]
    \node[rotate=270, scale=1.5, text opacity=0.5] 
        at ([xshift=-1cm, yshift=7cm]current page.south east) {\textbf{Ingeniería Matemática}};
\end{tikzpicture}

\begin{tikzpicture}[remember picture, overlay]
    \node[rotate=270, scale=1.5, text opacity=0.5] 
        at ([xshift=-1cm, yshift=16.1cm]current page.south east) {\textbf{Cálculo II}};
\end{tikzpicture}

\begin{tikzpicture}[remember picture, overlay]
    \draw[safetyorange(blazeorange), line width=2mm]
    (current page.south west) rectangle (current page.north west)
\end{tikzpicture}

}
%Carga el estilo
\fancypagestyle{Normal}{%
  \fancyhead[L]{\emph{Departamento de Matemáticas}}
  \fancyhead[R]{\emph{Univesidad de Santiago de Chile}}
  
  \fancyhead[C]{\Estilo}
  \renewcommand{\headrulewidth}{0pt}
}%



%Estilo de Pagina resumen
\newcommand{\mypage}{%
    \begin{tikzpicture}[remember picture, overlay]
        
        \fill[fill=safetyorange(blazeorange)](current page.north west) rectangle ([yshift=-22mm]current page.north east) node[midway] {\Huge\bfseries\sffamily\color{white}{Resumen}};

        
    \end{tikzpicture}
}%
%Carga el estilo
\fancypagestyle{Resumen}{%
  \fancyhf{}
  \fancyhead[C]{\mypage}
  \renewcommand{\headrulewidth}{0pt}
}%
\renewcommand{\footrulewidth}{0.5pt}

 \NewTcbTheorem[number within=section]{mytheo}{My Theorem}%
 {colback=green!5,colframe=green!35!black,fonttitle=\bfseries}{th}

%Ejemplo
\NewTcbTheorem[number within=section]{eje}{Ejemplo}
{ blanker,breakable,left=5mm,
 before skip=10pt,after skip=10pt,
 borderline west={1mm}{0pt}{red}}

%Problemas comunes
\NewTcbTheorem[number within=section]{prob}{Problema}{colback=gray!15!white
,title= Definition, attach boxed title to top left, fonttitle=\bfseries, sharp corners,boxrule=0pt,coltitle=black, boxed title style={colback=black!10}}{theo}

 %Propuestos
 \tcbset{ribbonbox/.style={enhanced,colback=gray!15!white, sharpish corners,  frame style={left color=red!75!black,
 right color=blue!75!black}, overlay={\path[fill=blue!75!white,draw=blue,double=white!85!blue,
 preaction={opacity=0.6,fill=blue!75!white},
 line width=0.1mm,double distance=0.2mm,
 pattern=fivepointed stars,pattern color=white!75!blue]
 ([xshift=-0.2mm,yshift=-1.02cm]frame.north east)-- ++(-1,1)-- ++(-0.5,0)-- ++(1.5,-1.5)-- cycle;}}}

 %DiseñoTCB
\NewTcbTheorem[number within=section]{defi}{Definición}{colback=bleudefrance!10,enhanced,title=Definition,
	attach boxed title to top left={xshift=-4mm},boxrule=0pt,after skip=1cm,before skip=1cm,right skip=0cm,breakable,fonttitle=\bfseries,toprule=0pt,bottomrule=0pt,rightrule=0pt,leftrule=4pt,arc=0mm,skin=enhancedlast jigsaw,sharp corners,colframe=bleudefrance,colbacktitle=bleudefrance, boxed title style={
		frame code={ 
			\fill[bleudefrance](frame.south west)--(frame.north west)--(frame.north east)--([xshift=3mm]frame.east)--(frame.south east)--cycle;
			\draw[line width=1mm,bleudefrance]([xshift=2mm]frame.north east)--([xshift=5mm]frame.east)--([xshift=2mm]frame.south east);
			
			\draw[line width=1mm,bleudefrance]([xshift=5mm]frame.north east)--([xshift=8mm]frame.east)--([xshift=5mm]frame.south east);
			\fill[bleudefrance!40](frame.south west)--+(4mm,-2mm)--+(4mm,2mm)--cycle;}}}{th}

 \NewTcbTheorem[number within=section]{mytheo}{My Theorem}%
 {colback=bleudefrance!10,enhanced, title=Definition, attach boxed title to top left={xshift=-4mm}, boxrule=0pt, after skip=1cm,right skip=0cm, breakable, fonttitle=\bfseries, toprule= 0pt, bottomrule=0pt, rightrule=0pt, leftrule=0pt, arc=0mm, skin=enhancedlast jigsaw, sharpcorners, colframe=bleudefrance, colbacktitle=bleudefrance, boxed title style={
 frame code={ 
			\fill[bleudefrance](frame.south west)--(frame.north west)--(frame.north east)--([xshift=3mm]frame.east)--(frame.south east)--cycle;
			\draw[line width=1mm,bleudefrance]([xshift=2mm]frame.north east)--([xshift=5mm]frame.east)--([xshift=2mm]frame.south east);
			
			\draw[line width=1mm,bleudefrance]([xshift=5mm]frame.north east)--([xshift=8mm]frame.east)--([xshift=5mm]frame.south east);
			\fill[bleudefrance!40](frame.south west)--+(4mm,-2mm)--+(4mm,2mm)--cycle;}}}{th}

\NewTcbTheorem[number within=section]{teo}{Teorema}{colback=cadmiumred!10,enhanced,title=Definition,
	attach boxed title to top left={xshift=-4mm},boxrule=0pt,after skip=1cm,before skip=1cm,right skip=0cm,breakable,fonttitle=\bfseries,toprule=0pt,bottomrule=0pt,rightrule=0pt,leftrule=4pt,arc=0mm,skin=enhancedlast jigsaw,sharp corners,colframe=cadmiumred,colbacktitle=cadmiumred, boxed title style={
		frame code={ 
			\fill[cadmiumred](frame.south west)--(frame.north west)--(frame.north east)--([xshift=3mm]frame.east)--(frame.south east)--cycle;
			\draw[line width=1mm,cadmiumred]([xshift=2mm]frame.north east)--([xshift=5mm]frame.east)--([xshift=2mm]frame.south east);
			
			\draw[line width=1mm,cadmiumred]([xshift=5mm]frame.north east)--([xshift=8mm]frame.east)--([xshift=5mm]frame.south east);
			\fill[cadmiumred!40](frame.south west)--+(4mm,-2mm)--+(4mm,2mm)--cycle;}}}{th}

            
\NewTcbTheorem[number within=section]{prop}{Proposición}{colback=gold(web)(golden)!10,enhanced,title=Definition,
	attach boxed title to top left={xshift=-4mm},boxrule=0pt,after skip=1cm,before skip=1cm,right skip=0cm,breakable,fonttitle=\bfseries,toprule=0pt,bottomrule=0pt,rightrule=0pt,leftrule=4pt,arc=0mm,skin=enhancedlast jigsaw,sharp corners,colframe=gold(web)(golden),colbacktitle=gold(web)(golden), boxed title style={
		frame code={ 
			\fill[gold(web)(golden)](frame.south west)--(frame.north west)--(frame.north east)--([xshift=3mm]frame.east)--(frame.south east)--cycle;
			\draw[line width=1mm,gold(web)(golden)]([xshift=2mm]frame.north east)--([xshift=5mm]frame.east)--([xshift=2mm]frame.south east);
			
			\draw[line width=1mm,gold(web)(golden)]([xshift=5mm]frame.north east)--([xshift=8mm]frame.east)--([xshift=5mm]frame.south east);
			\fill[gold(web)(golden)!40](frame.south west)--+(4mm,-2mm)--+(4mm,2mm)--cycle;}}}{th}


\NewTcbTheorem[number within=section]{coro}{Corolario}{colback=kellygreen!10,enhanced,title=Definition,
	attach boxed title to top left={xshift=-4mm},boxrule=0pt,after skip=1cm,before skip=1cm,right skip=0cm,breakable,fonttitle=\bfseries,toprule=0pt,bottomrule=0pt,rightrule=0pt,leftrule=4pt,arc=0mm,skin=enhancedlast jigsaw,sharp corners,colframe=kellygreen,colbacktitle=kellygreen, boxed title style={
		frame code={ 
			\fill[kellygreen](frame.south west)--(frame.north west)--(frame.north east)--([xshift=3mm]frame.east)--(frame.south east)--cycle;
			\draw[line width=1mm,kellygreen]([xshift=2mm]frame.north east)--([xshift=5mm]frame.east)--([xshift=2mm]frame.south east);
			
			\draw[line width=1mm,kellygreen]([xshift=5mm]frame.north east)--([xshift=8mm]frame.east)--([xshift=5mm]frame.south east);
			\fill[kellygreen!40](frame.south west)--+(4mm,-2mm)--+(4mm,2mm)--cycle;}}}{th}


 \pagestyle{Normal}
 %Sangía
\setlength{\parindent}{0pt}

 
\begin{document}


\chapter{Axiomas de Cuerpo}

\begin{prob}
\quad \quad
Sean $\forall a,b \in \mathbb{R}$, $a,b\neq 0$. Demuestre que $(ab)^{-1}=b^{-1}a^{-1}$ 
\end{prob}
\textbf{Solución:}
Para demostrar esta propiedad basta ver primero que $b^{-1}a^{-1}$ es el inverso multiplicativo de $ab$. En efecto:

\begin{align*}
    (ab)(b^{-1}a^{-1})&=a(bb^{-1})a^{-1} && \textbf{(Asociatividad de $\cdot$)} \\
    &= a \cdot 1 \cdot a^{-1} && \textbf{(Inverso Multiplicativo)} \\
    &= a \cdot a^{-1} && \textbf{(Neutro Multiplicativo)} \\
    &= 1 && \textbf{(Inverso Multiplicativo)}
\end{align*}

Para concluir, utilizamos el teorema {\Large Agregar cita}, el cual nos dice que el inverso multiplicativo es único. Así $(ab)^{-1}=b^{-1} a^{-1}$

\begin{prob}
    \quad \quad
Sean $a,b \in \mathbb{R}$, $a,b \neq 0$. Demuestre que $a^{-1}\cdot b= ( a \cdot b^{-1}\cdot)^{-1}$
\end{prob}
\textbf{Solución:}
Tenemos que verificar que $( a \cdot b^{-1})$ es el inverso multiplicativo de $a^{-1}\cdot b$. Notemos que 

\begin{align*}
    (a^{-1}\cdot b) \cdot (a \cdot b^{-1}) &=   (a^{-1}\cdot b) \cdot ( b^{-1}\cdot a ) && \textbf{(Conmutatividad)} \\
    &=  a^{-1} \cdot (b \cdot b^{-1}) \cdot a  && \textbf{(Asociatividad de $\cdot$)} \\
    &=  a^{-1} \cdot 1 \cdot a && \textbf{(Inverso Multiplicativo)} \\
    &= a^{-1} \cdot a && \textbf{(Neutro Multiplicativo)} \\
    &= 1 && \textbf{(Inverso Multiplicativo)}
\end{align*}

Ahora dado que el inverso multiplicativo es único, entonces tenemos que $a^{-1}\cdot b= (a \cdot b^{-1})^{-1}$.

\begin{prob}
    \quad \quad 
    Demuestre que $\forall a \in \mathbb{R}$ se tiene que $(-1)\cdot a = -a$
\end{prob}
\textbf{Solución:} 
Tenemos que probar que $a$ es el inverso aditivo de $(-1) \cdot a$. Luego 
\begin{align*}
    a + (-1)\cdot a  &= a \cdot 1 + (-1)\cdot a && \textbf{(Neutro Multiplicativo)} \\
    &= a \cdot (1+(-1)) && \textbf{(Distributividad)} \\
    &= a \cdot (0) && \textbf{(Inverso Aditivo)} \\
    &= 0 
\end{align*}
Utilizando el problema \ref{prob: 1.0.4} se tiene lo último

\begin{prob}
\quad \quad 
    Sea $a \in \mathbb{R}$. Pruebe que $a \cdot 0 = 0$.
\end{prob} \label{prob: 1.0.4}
\textbf{Solución:}

Notemos que 

\begin{align*}
   a \cdot 0 &= a \cdot (0 + 0 )  && \textbf{(Neutro de la Suma)} \\
 a\cdot 0 &= a \cdot 0 + a \cdot 0 && \textbf{(Distributividad)} \\
  a\cdot 0 + (-(a\cdot 0)) &= a\cdot 0 + a \cdot 0 + (-(a\cdot 0)) && \textbf{(Exist. Inverso)} \\
  0&= a\cdot 0 + (  a \cdot 0 -(a\cdot 0) ) && \textbf{(Asociatividad)} \\
  0&= a\cdot 0 && \textbf{(Inverso Aditivo)}
\end{align*}

\begin{prob}
\quad \quad 
    Demuestre que $\forall a, b \in \mathbb{R}$ se tiene que $-(a+b) = (-a) + (-b)$ .
\end{prob}
\textbf{Solución:} Notemos que

\begin{align*}
    (a+b)+ [(-a)+(-b)] &= a + [b+ (-a)] + (-b)   && \textbf{(Asociatividad)} \\
    &= a + [(-a) + b ]+ (-b) && \textbf{(Conmutatividad)} \\
    &= [a + (-a)]+ [b + (-b)] && \textbf{(Asociatividad)} \\
    &= 0 + 0  && \textbf{(Inverso Aditivo)} \\
    &=0 && \textbf{(Neutro Aditivo)}
\end{align*}

\begin{prob}
\quad \quad 
    Demuestre que $\forall a, b \in \mathbb{R}$ es tienen las siguientes propiedades 

    \begin{itemize}
        \item[\textbf{a)}] $-(-a)=a$
        \item[\textbf{b)}] $(-a) \cdot b=-(ab)$ 
        \item[\textbf{c)}] $a \cdot (-b) = -(ab)$
        \item[\textbf{d)}] $(-a)(-b)=ab$
    \end{itemize}
\end{prob}
\textbf{Solución:}
\begin{itemize}
    \item[\textbf{a)}] No es difícil notar que 
    \begin{align*}
        a+ (-a)&=0 \\
        (-a)+a&=0
    \end{align*}
    Esto nos dice que $a$ es el inverso aditivo de $-a$. Por lo tanto, dado que el inverso aditivo es único, se tiene que $a = -(-a)$.
      \item[\textbf{b)}] En efecto

      \begin{align*}
          (-a) \cdot b + a \cdot b&=((-a) + a ) \cdot b &&\textbf{(Distributividad)} \\
          &= 0 \cdot b && \textbf{(Inverso Aditivo)} \\
          &=0
      \end{align*}
\item[\textbf{c)}] Análogamente al ejercicio anterior 

\begin{align*}
    a \cdot (-b) + a\cdot b &= a \cdot (b + (-b)) && \textbf{(Distributividad)} \\
    &=a \cdot 0 && \textbf{(Inverso Aditivo)} \\
    &=0
\end{align*}
      \item[\textbf{d)}] Recordemos la propiedad $(-a)=(-1) \cdot a$. Entonces tenemos que 

      \begin{align*}
          (-a)(-b)&= [(-1)a][(-1)b]  \\
          &= (-1)\cdot[a \cdot (-1)] \cdot b && \textbf{(Asociatividad)} \\
          &=(-1) \cdot [(-1) \cdot a] \cdot b && \textbf{(Conmutatividad)} \\
          &= [(-1)(-1)]\cdot [ab] && \textbf{(Asociatividad)} \\
          &=1 \cdot ab 
          &=ab
      \end{align*}

      {\Large} Aquí falta probar lo $(-1)(-1)=1$ para que esté listo.
\end{itemize}
\begin{prob}
\quad \quad 
    Demuestre que $\forall a, b \in \mathbb{R}$, $a,b \neq 0$ se tienen las siguientes propiedades

    \begin{itemize}
        \item[\textbf{a)}] $(a^{-1})^{-1}=a$
        \item[\textbf{b)}] $a^{-1} \cdot b = ( a \cdot b^{-1})^{-1}$
        \item[\textbf{c)}] $a \cdot b^{-1} = (a^{-1} \cdot b)^{-1}$ 
        \item[\textbf{d)}] $(ab)^{-1}=a^{-1} \cdot b^{-1}$
    \end{itemize}
\end{prob}
\textbf{Solución:}

\begin{itemize}
    \item[\textbf{a)}]   No es difícil notar que 
\begin{align*}
    a \cdot a^{-1}&=1 \\
    a^{-1} \cdot a&=1
\end{align*}
Esto quiere decir que $a^{-1}$ es el inverso multiplicativo de $a$. Ahora dado que el inverso multiplicativo es único, se tiene que  $(a^{-1})^{-1}=a$.
\item[\textbf{b)}] Tenemos que verificar que $( a \cdot b^{-1})$ es el inverso multiplicativo de $a^{-1}\cdot b$. Notemos que 

\begin{align*}
    (a^{-1}\cdot b) \cdot (a \cdot b^{-1}) &=   (a^{-1}\cdot b) \cdot ( b^{-1}\cdot a ) && \textbf{(Conmutatividad)} \\
    &=  a^{-1} \cdot (b \cdot b^{-1}) \cdot a  && \textbf{(Asociatividad de $\cdot$)} \\
    &=  a^{-1} \cdot 1 \cdot a && \textbf{(Inverso Multiplicativo)} \\
    &= a^{-1} \cdot a && \textbf{(Neutro Multiplicativo)} \\
    &= 1 && \textbf{(Inverso Multiplicativo)}
\end{align*}

Ahora dado que el inverso multiplicativo es único, entonces tenemos que $a^{-1}\cdot b= (a \cdot b^{-1})^{-1}$.
\item[\textbf{c)}] Análogamente al problema \textbf{b)} tenemos que 

\begin{align*}
    (a\cdot b^{-1}) \cdot (a^{-1}\cdot b) &= (a\cdot b^{-1}) \cdot (b \cdot a^{-1}) && \textbf{(Conmutatividad)} \\
    &= a \cdot ( b^{-1} \cdot b) \cdot a^{-1} && \textbf{(Asociatividad)} \\
    &= a \cdot 1 \cdot a^{-1} && \textbf{(Inverso Multiplicativo)} \\
    &= a \cdot a^{-1} && \textbf{(Neutro Multiplicativo)} \\
    &= 1 && \textbf{(Inverso Multiplicativo)}
\end{align*}

\item[\textbf{d)}] Para demostrar esta propiedad basta ver primero que $a^{-1}b^{-1}$ es el inverso multiplicativo de $ab$. En efecto:

\begin{align*}
    (ab)(a^{-1}b^{-1})&=(ab)(b^{-1}a^{-1}) && \textbf{(Conmutatividad)} \\
    &=a(bb^{-1})a^{-1} && \textbf{(Asociatividad de $\cdot$)} \\
    &= a \cdot 1 \cdot a^{-1} && \textbf{(Inverso Multiplicativo)} \\
    &= a \cdot a^{-1} && \textbf{(Neutro Multiplicativo)} \\
    &= 1 && \textbf{(Inverso Multiplicativo)}
\end{align*}

Para concluir, utilizamos el teorema {\Large Agregar cita}, el cual nos dice que el inverso multiplicativo es único. Así $(ab)^{-1}=a^{-1} b^{-1}$
\end{itemize}
\begin{prob}{}{}
 Demuestre que para $\forall a,b \in \mathbb{R
 }$, $a,b \neq 0$, se tiene que 
\begin{align*}
     -(a^{-1})=(-a)^{-1}
\end{align*}
\end{prob}
\textbf{Solución:} No es difícil notar que existen dos vías para resolver este problema. Una es utilizando los inversos aditivos, y el otro, los inversos multiplicativos. En este caso, elegiremos la segunda vía. Notamos que 

\begin{align*}
    -(a^{-1})\cdot (-a)&=[(-1)a^{-1}] \cdot [(-1)a] && \\
    &= [(-1)a^{-1}] \cdot [a \cdot (-1)] && \textbf{(Conmutatividad)} \\
    &= (-1) \cdot [a^{-1} \cdot a] \cdot (-1) && \textbf{(Asociatividad)} \\
    &= (-1) \cdot 1 \cdot (-1) && \textbf{(Inverso Multiplicativo)} \\
    &=(-1)(-1) && \textbf{(Neutro Multiplicativo)} \\
    &=1
\end{align*}

\begin{prob}{}{}
    Sean $a,b,c \in \mathbb{R}$, $b,c\neq 0$. Pruebe que 
    \begin{align*}
        (ac^{-1}+b^{-1})[(bc)\cdot (ab+c)^{-1}]=1
    \end{align*}
\end{prob}
\textbf{Solución:}

Notemos que 
\begin{align*}
    (ac^{-1}+b^{-1})[(bc)\cdot (ab+c)^{-1}] &=  [(ac^{-1}+b^{-1})\cdot(bc)]\cdot (ab+c)^{-1} && \textbf{(Asociatividad)} \\
    &=[(ac^{-1})(bc)+b^{-1}(bc)] \cdot (ab+c)^{-1} && \textbf{(Distributividad)} \\
    &=[(ac^{-1})(cb)+b^{-1}(bc)] \cdot (ab+c)^{-1} && \textbf{(Distributividad)} \\
    &=[a(c^{-1}c)b+(b^{-1}b)c] \cdot (ab+c)^{-1} && \textbf{(Asociatividad)} \\
    &=[a \cdot 1 \cdot b + 1 \cdot c ] \cdot (ab+c)^{-1} && \textbf{(Inverso Multiplicativo)} \\
    &=(ab+c) \cdot (ab+c)^{-1}   && \textbf{(Neutro Multiplicativo)} \\
    &=1
\end{align*}

\begin{prob}
    Sean $a,b \in \mathbb{R}$, $a,b \neq 0$. Demuestre que $(a+b)(ab)^{-1}=a^{-1}+b^{-1}$
\end{prob}
\textbf{Solución:}
En efecto

\begin{align*}
    (a+b)(ab)^{-1}=a^{-1}+b^{-1}&=(a+b)(a^{-1}b^{-1}) && \textbf{Propiedad (*)} \\
    &=a\cdot (a^{-1} b^{-1}) + b \cdot (a^{-1}b^{-1}) && \textbf{(Distributividad)} \\
    &= (a a^{-1})b^{-1} +  b \cdot (a^{-1}b^{-1}) && \textbf{(Asociatividad)} \\
    &=  (a a^{-1})b^{-1} +  b \cdot ( b^{-1}a^{-1}) && \textbf{(Conmutatividad)} \\
    &=  (a a^{-1})b^{-1} +  (bb^{-1})a^{-1} && \textbf{(Asociatividad)} \\
    &= 1\cdot b^{-1} +1 \cdot a^{-1} && \textbf{(Inverso Multiplicativo)} \\
    &=b^{-1} +a^{-1} && \textbf{(Neutro Multiplicativo)} \\
    &= a^{-1}+b^{-1} && \textbf{(Conmutatividad)} 
\end{align*}

\begin{prob}
    \quad \quad
    Utilizando la notación $a \cdot a = a^2$. Demuestre que $\forall a,b \in \mathbb{R}\setminus \{0\}$,
    \begin{align*}
        (a^{-1}+(-b)^{-1})\cdot (a^{-1}b)=(a^{-1})^2 b + (-a)^{-1}
    \end{align*}
\end{prob}
\textbf{Solución:}

 En efecto

 \begin{align*}
      (a^{-1}+(-b)^{-1})\cdot (a^{-1}b)&= a^{-1}(a^{-1}b)+(-b)^{-1}(a^{-1}b) && \textbf{(Distributividad)} \\
      &= (a^{-1}a^{-1})b+(-b)^{-1}(a^{-1}b) && \textbf{(Asociatividad)} \\
      &=(a^{-1}a^{-1})b+(-b)^{-1}(ba^{-1}) && \textbf{(Conmutatividad)} \\
      &=(a^{-1}a^{-1})b+(-(b^{-1}b)a^{-1}) && \textbf{(Asociatividad)} \\
      &=(a^{-1})^2b+ (-(a^{-1})) && \textbf{(Inverso Multiplicativo)} \\
      &=(a^{-1})^2b+(-a)^{-1} && \textbf{(Problema 1.0.8)} \\
 \end{align*}
\begin{prob}
     \quad \quad
    Utilizando la notación $a \cdot a = a^2$. Demuestre que $\forall a,b \in \mathbb{R}\setminus \{0\}$, 
    \begin{align*}
        (a^2 + (-(b^{-1})^2))\cdot ((a^{-1})^2 b )= b + (-b)^{-1}(a^{-1})^2
    \end{align*}
\end{prob}
\textbf{Solución:}
En efecto, notemos que 

\begin{align*}
   (a^2 + (-(b^{-1})^2))\cdot ((a^{-1})^2 b )   &=a^2 ((a^{-1})^2 b ) + (-(b^{-1})^2 (a^{-1})^2 b ) && \textbf{(Distributividad)} \\
   &=(a^2 (a^{-1})^2)b  + (-(b^{-1})^2 (a^{-1})^2 b ) && \textbf{(Asociatividad)} \\
   &= b + (-(b^{-1})^2 (a^{-1})^2 b ) && \textbf{(Inverso Multiplicativo)} \\
   &= b + (-(b^{-1})^2  b (a^{-1})^2  ) && \textbf{(Conmutatividad)} \\
   &=b + (-(b^{-1}(a^{-1})^2  )) && \textbf{(Inverso Multiplicativo)} \\
   &=b + (-(b^{-1}))(a^{-1})^2 && \textbf{(Asociatividad)} \\
   &=b+ (-b)^{-1}(a^{-1})^2
\end{align*}
\begin{prob}
\quad \quad
    Sea $a \in \mathbb{R}$ arbitrario. Si $b \in \mathbb{R}$ es tal que $a+b=a$, entonces $b=0$
\end{prob}
\textbf{Solución:}
Como $a \in \mathbb{R}$ por \textbf{(A4)} sabemos que existe su inverso aditivo $-a \in \mathbb{R}$, que cumple $a + (-a)=0$. Por lo tanto, sumemos en ambas partes de la igualdad $-a$. 

\begin{align*}
    (a+b)+(-a)&=a+(-a) \\
    (b+a)+(-a)&=0  && \textbf{(Conmutatividad)} \\
    b+ (a+(-a))&=0 && \textbf{(Asociatividad)} \\
    b+0&=0 && \textbf{(Inverso Aditivo)} \\
    b=0
\end{align*}

\begin{prob}
\quad \quad
    Sean $a,b,c \in \mathbb{R}$. Demuestre que se cumplen las siguentes propiedades:

    \begin{itemize}
        \item[\textbf{a)}] Si $a+b=a+c  \ \Longleftrightarrow\ b=c$
        \item[\textbf{b)}] Si $ab =ac$ $\Longrightarrow$ $b=c$
        \item[\textbf{c)}] Si $ab=0$ $\Longrightarrow$ $(a=0)$ $\vee$ $(b=0)$
    \end{itemize}
\end{prob}
\textbf{Solución:}
\begin{itemize}
    \item[\textbf{a)}] Probemos de izquierda a derecha. Notemos que $a \in \mathbb{R}$, entonces existe el inverso aditivo $-a \in \mathbb{R}$ tal que $a+(-a)=0$ por \textbf{(A4)}. Luego aplicando ambos lados de la igualdad.

    \begin{align*}
        (-a)+(a+b)&=(-a)+(a+c) \\
        ((-a)+a)+b&=((-a)+a)+c && \textbf{(Asociatividad)} \\
        0+b&=0+c && \textbf{(Inverso Aditivo)} \\
        b&=c && \textbf{(Neutro Aditivo)}
    \end{align*}
La otra implicancia es directa
    
    \item[\textbf{b)}] Dado que $a \neq 0$, entonces por \textbf{(A4)} sabemos que existe su inverso multiplicativo $a^{-1}$. Entonces multiplicando a ambos lados de la igualdad obtenemos

    \begin{align*}
        a^{-1}(ab)&=a^{-1}(ac) 
        (a^{-1}a)b&=(a^{-1})ac &&  \textbf{(Asociatividad)} \\
        1\cdot b &= 1 \cdot c && \textbf{(Inverso Multiplicativo)} \\
        b&=c && \textbf{(Neutro Multiplicativo)}
    \end{align*}
    
    \item[\textbf{c)}]  Supongamos que $a \neq 0$, ya que en el caso $a=0$, no hay nada que probar. Luego dado que $a,b \neq 0$, entonces existen los inversos multiplicativos $a^{-1}$ y $b^{-1}$. Análogamente a la parte \textbf{a)}, tenemos que

    \begin{itemize}
        \item[\textbf{Caso 1:}]   
        \begin{align*}
            a^{-1}(ab)=  a^{-1} \cdot 0 \quad \Longrightarrow \quad b=0
        \end{align*}

          \item[\textbf{Caso 2:}]

            \begin{align*}
            (ab)b^{-1}=0 \cdot b^{-1} \quad \Longrightarrow \quad a=0
          \end{align*}
    En cualquiera de los dos casos se cumple la proposición  $(a=0)$ $\vee$ $(b=0)$.
    \end{itemize}    
\end{itemize}





\begin{prob}
    Sean $a,b,c,d \in \mathbb{R}$, $b,d \neq 0$. Demuestre que se cumplenlas siguientes propiedades

    \begin{itemize}
        \item[\textbf{a)}] $ab^{-1} + cd^{-1}=(ad+cb)(bd)^{-1}$

        \item[\textbf{b)}] $(ab^{-1})\cdot (cd^{-1})=(ac)\cdot (bd)^{-1}$

        
    \end{itemize}
\end{prob}
\textbf{Solución:}

\begin{itemize} 
     \item[\textbf{a)}] Empezamos por la parte derecha de la igualdad

     \begin{align*}
         (ad+cb)(bd)^{-1}&=(ad+cb)(b^{-1}d^{-1}) && \textbf{(Problema 1.0.7, d))} \\
         &=(ad+cb)(d^{-1}b^{-1}) && \textbf{(Conmutatividad)} \\
        &=ad(d^{-1}b^{-1})+cb(d^{-1}b^{-1}) && \textbf{(Distributividad)} \\
        &=a(dd^{-1})b^{-1}++cb(d^{-1}b^{-1}) && \textbf{(Asociatividad)} \\
        &=a(dd^{-1})b^{-1}++cb(b^{-1}d^{-1}) && \textbf{(Conmutatividad)} \\
        &=a(dd^{-1})b^{-1}++c(bb^{-1})d^{-1} && \textbf{(Asociatividad)} \\
        &=a \cdot 1 \cdot b^{-1}+ c \cdot 1 \cdot d^{-1} && \textbf{(Inverso Multiplicativo)} \\
        &=ab^{-1}+cd^{-1} && \textbf{(Neutro Multiplicativo)}
     \end{align*}

      \item[\textbf{b)}] En efecto 

      \begin{align*}
          (ab^{-1})\cdot (cd^{-1})&= a (b^{-1}c) d^{-1} && \textbf{(Asociatividad)} \\
          &=a (c b^{-1})d^{-1} && \textbf{(Conmutatividad)} \\
          &= (ac) \cdot (b^{-1}d^{-1}) && \textbf{(Asociatividad)}\\
          &= (ac) \cdot (bd)^{-1} && \textbf{(Problema 1.0.7, d))}\\
      \end{align*}
\end{itemize}

\begin{mytheo}
    Hola
\end{mytheo}

\end{document}
